\subsection{Overview}

A feasibility study evaluates whether the proposed system can be successfully developed within the available resources, time, technical knowledge, and budget. The objective is to determine whether Dikaios Amanah is practical as a university-level Software Engineering capstone project while also demonstrating the potential for future real-world deployment.

The proposed system has been designed as a Minimum Viable Product (MVP), ensuring that essential features can be implemented within the academic timeline while allowing future expansion beyond the scope of this project.

%----------------------------------------------------

\subsection{Technical Feasibility}

The proposed system utilizes mature and widely adopted web development technologies that are well supported by extensive documentation and open-source communities.

The planned technology stack consists of:

\begin{table}[H]
\centering
\begin{tabular}{|p{5cm}|p{8cm}|}
\hline
\textbf{Technology} & \textbf{Purpose} \\
\hline
React.js & Frontend development \\
\hline
Node.js & Backend runtime environment \\
\hline
Express.js & REST API development \\
\hline
PostgreSQL & Relational database management \\
\hline
Firebase Storage & Secure document storage \\
\hline
JWT & User authentication \\
\hline
bcrypt & Password hashing \\
\hline
SHA-256 & Document integrity verification \\
\hline
Git \& GitHub & Version control \\
\hline
Permissioned Trust and Audit Layer & Tamper-evident audit records \\
\hline
\end{tabular}
\caption{Proposed Technology Stack}
\end{table}

These technologies are suitable for building secure, scalable, and maintainable web applications. Their widespread adoption also ensures long-term maintainability and extensive learning resources.

The proposed layered architecture separates presentation, business logic, security, and data management, reducing development complexity and improving software maintainability.

Therefore, the project is considered technically feasible.

%----------------------------------------------------

\subsection{Economic Feasibility}

The project is intended for academic demonstration and therefore minimizes operational costs by utilizing free and open-source software wherever possible.

The estimated development cost is limited primarily to internet access, optional cloud hosting, domain registration, and testing resources.

Open-source technologies such as React, Node.js, PostgreSQL, Express.js, and GitHub significantly reduce software licensing costs.

Cloud services such as Firebase Storage also provide generous free usage tiers suitable for MVP development.

Consequently, the project can be completed with minimal financial investment.

%----------------------------------------------------

\subsection{Operational Feasibility}

The proposed platform has been designed with usability and accessibility in mind.

Business owners will be able to submit investment campaigns through an intuitive dashboard, while investors will have access to verified business information before making investment decisions.

System administrators will manage verification, approvals, and platform monitoring through a dedicated administrative interface.

Role-Based Access Control (RBAC), secure authentication, activity logging, and document verification improve operational security while simplifying day-to-day platform management.

The system architecture supports future scalability without requiring major architectural changes.

Therefore, the project is considered operationally feasible.

%----------------------------------------------------

\subsection{Schedule Feasibility}

The proposed project has been divided into manageable development phases following an iterative Software Development Life Cycle (SDLC).

Major development phases include:

\begin{itemize}

\item Requirements Analysis

\item System Design

\item Database Design

\item UI/UX Design

\item Backend Development

\item Frontend Development

\item Security Implementation

\item Trust and Audit Layer Development

\item Testing and Debugging

\item Documentation

\end{itemize}

The MVP intentionally excludes complex production-level features such as live banking integration, regulatory compliance, and smart contract automation to ensure successful completion within the academic schedule.

Based on the planned scope and available development timeline, the project is considered schedule feasible.

%----------------------------------------------------

\subsection{Risk Assessment}

Although the project is considered feasible, several risks have been identified.

\begin{table}[H]
\centering
\begin{tabular}{|p{6cm}|p{3cm}|p{4cm}|}
\hline
\textbf{Risk} & \textbf{Impact} & \textbf{Mitigation Strategy} \\
\hline
Development Time Constraints & Medium & Prioritize MVP features. \\
\hline
Learning New Technologies & Medium & Incremental learning and prototyping. \\
\hline
Authentication Security Issues & High & Follow secure coding practices and testing. \\
\hline
Database Design Errors & Medium & Early ERD design and normalization. \\
\hline
Document Verification Challenges & Medium & Administrative review process. \\
\hline
Blockchain Integration Complexity & Medium & Implement a lightweight permissioned audit layer for the MVP. \\
\hline
\end{tabular}
\caption{Project Risk Assessment}
\end{table}

Proper planning, iterative development, version control, and continuous testing will significantly reduce these risks throughout the software development lifecycle.

%----------------------------------------------------

\subsection{Feasibility Conclusion}

Based on the technical, economic, operational, and schedule analyses, Dikaios Amanah is considered a feasible Software Engineering capstone project. The selected technology stack, clearly defined MVP scope, and phased development strategy provide a realistic pathway for successful implementation within the available academic timeframe while allowing future enhancements for production deployment.